% !TeX root = ../papers/nc-p3.tex
\section{正則正規な一次直線列を持つ窓と楕円三次商}
\label{sec:regular-normal-linear}
量子点への退化とは別の切り口として、一次の正規元を用いる。
正則正規な直線列とは、各 $i$ の非零元 $t_i\in C_{i,i+1}$ が、全 $i\le j$ で
\[
 C_{ij}t_j=t_i C_{i+1,j+1}
\]
を満たし、左右の乗法が単射であるような直線 $kt_i$ の列とする。周期は仮定しない。

\begin{theorem}\label{thm:regular-normal-linear}
この直線列を持つ標準14641型の添字代数は全て最小周期1である。
直線列付きの全同型類は、三生成三関係の標準Koszul左右AS正則代数
$S=T(V)/(R)$ と二写像 $\alpha:R\to V$, $\beta:R\to k$ から次のように与えられる。
$0\ne w\in(R\otimes V)\cap(V\otimes R)$ を取り、
\[
 h=(\alpha\otimes1-1\otimes\alpha)w,\qquad
 b=(\beta\otimes1-1\otimes\beta)w.
\]
採用する条件は $h\in R$, $\alpha(h)=-b$, $\beta(h)=0$ の三つである。
中心文字 $t$ を加え、$r=t\alpha(r)+t^2\beta(r)$ ($r\in R$) を課して得る代数を $A$ とする。
対象は $\check A$ と $kt$ である。二組の直線列付き同型は、
$(M\otimes M)R=R'$ となる可逆写像 $M:V\to V'$ と $e\in V^*$ が存在し、
全 $r\in R$, $r'=(M\otimes M)r$ について
\begin{align*}
 M\alpha(r)&=\alpha'(r')+(e\otimes M+M\otimes e)(r),\\
 \beta(r)+e(\alpha(r))&=\beta'(r')+(e\otimes e)(r)
\end{align*}
となることと同値である。
\end{theorem}
\begin{proof}
$t_i$ の代表を選ぶ。正規性と単射性により $a t_j=t_i\phi(a)$ で定まる写像
$\phi:C_{ij}\to C_{i+1,j+1}$ は全単射である。積 $ab$ にこの等式を二回適用すれば
$\phi(ab)=\phi(a)\phi(b)$、また $\phi(t_i)=t_{i+1}$ である。
従って標準四移動との可換性を先に仮定せず、$\phi$ の反復で通常の二次代数 $A$ の
反復表示が得られる。$t$ は中心で正則となる。
$S=A/(t)$ は三生成二次代数で、中心元の短完全列からHilbert級数 $(1-z)^{-3}$ を持つ。
Koszul性は変換環のTorスペクトル系列で従う。$\operatorname{Tor}^A_q(k,S)$ は
$q=0,1$ にだけあり、それぞれ $k,k(-1)$ である。もし $\operatorname{Tor}^S_p(k,k)$
に最初の非線型次数があれば、下段のその成分に入る微分はなく、出る可能性のある
$d_2$ は $p-2$ 列の線型な上段に入るのでそこにも消えない。$A$ のKoszul性に反する。
二次双対の級数は $(1+z)^3$ となり、大域次元は3である。
中心正則元に対する変換環Ext系列
（$\operatorname{Ext}^q_A(S,A)$ は $q=1$ の $S(1)$ だけ）を両側で使えばAS双対性も降りる。

$V$ を $A_1$ に持ち上げる。$A_2$ の次元から、中心三関係以外は $r$ の三つの持上げで尽き、
右辺には $tV$ と $kt^2$ の係数だけが現れる。$A$ は $U=A/(t-1)$ の
語長フィルトレーションのRees代数で、$\operatorname{gr} U=S$ である。
実際 $a\in A_n$ を $a t^{-n}$ として $A[t^{-1}]_0$ に入れると、単射な次数片の
増加和が $U$ となり、その逐次商が $S_n$ である。逆向きには
$U=T(V)/(r-\alpha(r)-\beta(r))$ がPBW、すなわち $S\to\operatorname{gr} U$ が同型なら、
上の斉次表示はそのRees代数である。各 $n$ で $S_n$ の語基底とその低次の $t$ 倍が
斉次表示を張り、Rees代数で独立だからである。
Koszul基底に対するBraverman--Gaitsgoryの二次PBW定理によれば、PBW性は
三次重なり上の表示した三条件に同値である。符号は $r-\alpha(r)-\beta(r)$ の規約による。
交叉は $S$ の三次双対成分の双対で一次元なので $w$ 一つで足りる。
ここで要求するPBW性は、三生成商が非可換でも成り立つ随伴次数の等式である。

PBWの場合は再び $t$ が正則である。上記Tor系列を逆に使うと、$S$ がKoszulなので
両段は内部次数と全ホモロジー次数が等しく、$A$ もKoszulである。
Hilbert級数とExt系列から両側AS、大域次元4、標準型を得る。
最後に、直線を保つ添字同型では頂点別スカラーを選んで全ての $t_i$ の像を $t'_i$ とできる。
中心性と正則性からこの同型は $\phi$ と可換し、位置によらない次数付き同型となる。
一次の持上げは $x\mapsto Mx+e(x)t'$ である。関係に代入し、$t'V'$ と $kt'^2$ の
二つの独立な係数を比較すると定理の式となる。逆もこの代入と可逆性による。
\end{proof}

\begin{theorem}\label{thm:elliptic-central-line}
上記の三生成商の二文字点対応が、滑らかな平面三次曲線 $E$ 上の自己同型
$\sigma$ のグラフであるとする。このとき正則正規一次直線列は一意であり、
定理\ref{thm:regular-normal-linear}の判定は直線の指定を忘れた添字保持同型にも必要十分である。
二つの商の間で試すべき射影行列は、二曲線の変曲点から有限個に絞れる。
\end{theorem}
\begin{proof}
$L=\mathcal O_E(1)$ とする。$\sigma^*L\simeq L$ なら、$\sigma$ は平面の射影変換に延びる。
このとき各二文字関係をその平面グラフ上で評価すると、$E$ 上で消える平面二次式である。
三次曲線のイデアルには非零二次式がないので、全てが平面グラフ上で零になる。
点対応の第一射影が $E$ という仮定に反する。

もう一つの直線列の生成元を商へ写して $v_i\in S_1$ とする。
$S_1 v_{i+1}=v_i S_1$ である。一方が零なら、もう一方も零である。
実際グラフ上での一次評価は単射で、$L$ と $\sigma^*L$ の全切断は基点を持たない。
全て非零なら、同じ等式を $E$ 上で評価して、
二つの積の線型系の共通零因子は $\sigma^*\operatorname{div}(v_{i+1})$ と
$\operatorname{div}(v_i)$ になる。各残り因子の線型系に基点がないため
零因子は重複度を込めて等しい。従って $\sigma^*L\simeq L$ となり、上と矛盾する。
よって全て零であり、元の直線列は各次数で $kt$ そのものである。
商の次数付き同型は滑らかな三次曲線間の射影同型を誘導する。
九つの変曲点から射影枠を選び、その像の順序付き候補を試して $R$ の像も検査すればよい。
実際、変曲点中に非共線な三点があり、その三辺上には合計で高々六点しかない。
\end{proof}

\begin{corollary}\label{cor:elliptic-central-linear-slice}
任意の滑らかな楕円三次商 $S$ を固定する。$e\mapsto D(e)$,
$D(e)(r)=(e\otimes1+1\otimes e)r$ の像に対して
$\operatorname{Hom}(R,V)$ の線型補空間を取り、$\alpha$ をこの補空間に限ることができる。
そこで $h(\alpha)\in R$ を課した各 $\alpha$ に対し、
\[
 \big(b(\beta),\beta(h(\alpha))\big)=\big(-\alpha(h(\alpha)),0\big)
\]
という線型系の全解が、採用する定数欄 $\beta$ を与える。異なる解の同型は
前定理の有限個の射影行列と重み $(1,2)$ の一斉倍率で尽くされる。
従って原点以外では、重み付き射影空間内の重み1、2、3の方程式と有限群の作用で全特殊値を扱える。
\end{corollary}
\begin{proof}
$D(e)=0$, $e\ne0$ と仮定する。全 $r$ で $r(e,e)=0$ なので $([e],[e])$ は
二文字点対応上にある。平面の対角に制限した式 $r(x,x)$ は全て $e$ で一次部分が零だから、
対角とのスキーム交叉のこの点の接次元は2である。滑らかな曲線のグラフの閉部分スキームで
接次元が1を超えることはできない。従って $D$ は単射であり、前定理の持上げ変更で
各 $\alpha$ をただ一回だけ補空間に入れられる。残る二条件は $\alpha$ を固定すると
まさに示した $\beta$ の線型系となる。
\end{proof}

\begin{corollary}\label{cor:sklyanin-central-classification}
滑らかな三次曲線を点対応に持つ通常の三生成Sklyanin商を固定し、
巡回関係を $r_a=p x_bx_c+q x_cx_b+s x_a^2$ と書く（$s\ne0$）。
この商に属する全窓の同型類は一点と重み付き射影空間
$\mathbb P(1,1,1,2,2,2)$ の有限群商である。ここでは正則正規一次直線列を通る商を指定する。
その直線列は一意なので、得る類の一般次元5は何も指定しない添字同型類の次元でもある。
全特殊値が正則であり、全て最小周期1、いずれも同じ $\check{S[t]}$ に平坦特殊化する。
\end{corollary}
\begin{proof}
$w=\sum_a x_a\otimes r_a=\sum_a r_a\otimes x_a$ を取れる。
$h=\sum_a(\alpha(r_a)\otimes x_a-x_a\otimes\alpha(r_a))$ は交代テンソルである。
$s\ne0$ より $R$ の基底は三つの純平方係数で区別でき、$R\cap\bigwedge^2 V=0$ である。
従ってPBW条件は $h=0$ に等しい。すなわち $\alpha(r_a)=\sum_j m_{aj}x_j$ の
$3\times3$ 係数行列が対称であればよく、三つの $\beta(r_a)$ は全て自由である。
一次の持上げの変更 $e$ は対角三係数にそれぞれ $2s e(x_a)$ を加減する。
従ってちょうど一回の変更で全対角を消せ、非対角三係数と定数三係数が残る。
文字を一斉にスカラー倍するとこれらの重みは1と2である（倍率の逆を選んでも同じ商となる）。
$\operatorname{Aut}_{\mathrm{gr}}(S)$ のスカラーを除いた群は前定理により有限である。
この群の作用は各射影行列の持上げを定理の二式に代入し、再び対角を消せば得られ、
持上げの違いは既に除いた一斉倍率だけである。原点以外の倍率軌道は上記重み付き
射影空間の点に対応する。原点への特殊化は $(\alpha,\beta)\mapsto
(u\alpha,u^2\beta)$ で与えられる。各次数で商の基底の持上げの $t$ 倍が族の基底となるため平坦である。
\end{proof}

このSklyanin商では全次数正則性を三次PBW定理によらず確かめることもできる。
$L_a=t x_a-x_a t$,
$P_a=r_a-(t\alpha(r_a)+\alpha(r_a)t)/2-\beta(r_a)t^2$ とする。
$\sum(L_a\otimes P_a+P_a\otimes L_a)$ は一位置回転で符号を変える。
基底部分では $w$ の巡回性を使う。$m$ が対称であるため、線型補正は
$v\in V$ に対する $(tv-vt)\otimes(tv+vt)+(tv+vt)\otimes(tv-vt)$ の線型結合であり、
定数補正も同様に符号付き巡回和となる。六関係は $R$ への射影と三中心関係から独立なので
両隣接階数は6である。$t$ だけに正の余重みを与えて全体倍率を除くと $\check{S[t]}$ に流れる。
命題\ref{prop:unrestricted-good-reduction}が任意の $m,\beta$ に適用できる。
